\chapter{Phương pháp đề xuất}
\label{chap:method}
\section{Tổng quan hướng tiếp cận}
Với mỗi segment \(t\) gồm \(L\) token, Decoder Qwen2.5-3B-Instruct tạo các trạng thái ẩn và sử dụng trạng thái của token cuối cùng \(h_{t,L}\) làm đầu vào cho Decision Head. Thông qua phép chiếu tuyến tính và hàm sigmoid, Decision Head dự đoán điểm quyết định truy hồi:

\begin{equation}d_t = \sigma(w^\top h_{t,L} + b).\label{eq:gate}\end{equation}

Trong giai đoạn suy luận (inference), tại prefill và biên câu đủ điều kiện, hệ thống thực hiện truy hồi nếu truy vấn hợp lệ và \(d_t \geq \tau\). Khi truy hồi được thực hiện, điểm \(d_t\) được liên kết với quan sát mới và sử dụng làm hệ số Attention Mixing \(d\) trong các lần forward Decoder tiếp theo, cho đến khi quan sát đó được thay thế bởi kết quả truy hồi mới.

Đối với các thực nghiệm kiểm chứng giả thuyết H1/H2, hệ số \(d\) được thiết lập trực tiếp thay vì dự đoán thông qua Decision Head, nhằm khảo sát ảnh hưởng của Attention Mixing đến LM Loss và gradient.

\section{Dữ liệu}
\textbf{Dữ liệu kiểm chứng giả thuyết H1/H2.} Các thực nghiệm sử dụng mẫu Multi-hop QA từ MuSiQue với dữ liệu được tạo bởi teacher model, phân nhóm theo số bước suy luận và tỷ lệ tài liệu hỗ trợ so với tài liệu gây nhiễu. Trong mỗi điều kiện thực nghiệm, câu hỏi, tài liệu và đích dự đoán được giữ cố định, trong khi hệ số Attention Mixing được thay đổi có kiểm soát nhằm khảo sát sự biến thiên của LM Loss và gradient.

\textbf{Dữ liệu huấn luyện ART.} Dữ liệu huấn luyện được xây dựng từ tập huấn luyện MuSiQue. Teacher model Qwen3-8B tạo chuỗi đích ngoại tuyến với bốn ví dụ few-shot và hướng dẫn xử lý trường hợp bằng chứng chưa đầy đủ. Trong Collect–Replay, ART dùng chuỗi đích này để thu thập quỹ đạo truy hồi rồi tối ưu mô hình trên các tiền tố đã lưu.

\textbf{Dữ liệu đánh giá.} Hệ thống được đánh giá trên ba bộ dữ liệu Multi-hop QA gồm MuSiQue, HotpotQA và 2WikiMultiHopQA. Các mẫu đánh giá được tách khỏi dữ liệu huấn luyện và được sử dụng để so sánh ART với các baseline No-RAG và Vanilla RAG\cite{musique,hotpot,twowiki}.

\section{Phương pháp}
\subsection{Attention Mixing}
Từ cùng dòng trạng thái, hệ thống chiếu \(Q,K,V\), rồi tính hai biến thể RoPE cho truy vấn và khóa. Đặt \(C\) là mặt nạ nhân quả, \(D\) là mặt nạ khóa của đoạn quan sát mới nhất và \(m\) là chiều khóa. Hai phân bố là
\begin{align}
A_+ &= \operatorname{softmax}(Q_+K_+^\top/\sqrt{m}+C),\\
A_- &= \operatorname{softmax}(Q_-K_-^\top/\sqrt{m}+C+D),\\
O &= \operatorname{Dropout}\big(dA_++(1-d)A_-\big)V_+.
\label{eq:mix}
\end{align}
Nhánh trừ hiệu chỉnh vị trí sau đoạn che bằng độ dài đoạn. Hai nhánh dùng chung giá trị của dòng có tài liệu. Sau phép trộn là phép chiếu đầu ra. Nếu thiếu mặt nạ hoặc thiếu \(d\), dùng nhánh cộng. Với các phép chiếu cố định, phần khối lượng trực tiếp trên khóa bị che tỉ lệ với \(d\); sau nhiều lớp trạng thái trung gian cũng phụ thuộc \(d\). Nhánh trừ vẫn có thể chịu ảnh hưởng của quan sát cũ trong tiền tố.
\Cref{fig:attention} minh họa hai nhánh Attention Mixing dùng chung value. Cả hai nhận cùng trạng thái đầu vào.
\par\medskip\noindent
\begin{minipage}{\linewidth}
\captionsetup{hypcap=false}
\centering
\includegraphics[width=.72\linewidth]{graphics/attention-figure.pdf}
\captionof{figure}{Trộn hai phân bố chú ý từ cùng trạng thái. Nhánh trừ che quan sát mới nhất, trong khi hai nhánh dùng chung giá trị. Sơ đồ do đồ án dựng từ mã hiện thực.}
\label{fig:attention}
\end{minipage}\par\medskip
\subsection{Mục tiêu học và giả thuyết}
Mục tiêu của ART là học khi nào cần truy hồi từ sự đánh đổi giữa khả năng dự đoán chuỗi đích và một số hạng phạt trên điểm quyết định. Với các vị trí token đích hợp lệ \(\mathcal T\), hàm mục tiêu trong Replay có dạng
\begin{equation}
\mathcal{L}_{\mathrm{LM}}=-\frac{1}{|\mathcal{T}|}
\sum_{t\in\mathcal{T}}\log p_\theta(y_t\mid x_{<t}),\qquad
\mathcal{L}=\mathcal{L}_{\mathrm{LM}}+0.05\sum_{j\in J}d_j.
\label{eq:loss}
\end{equation}
Token câu hỏi và tài liệu truy hồi không thuộc \(\mathcal T\). LM Loss dùng chuỗi đích do teacher model sinh. \(J\) là đa tập các lượt cộng điểm trong Replay, kể cả khoản cộng ở nhánh kết thúc nếu được kích hoạt. Regularization phạt điểm cao; LM Loss tạo tín hiệu theo mức độ mà quan sát cải thiện dự đoán. Số hạng phạt không đếm số truy vấn.

Khi Replay, gradient từ dự đoán token đi qua Attention Mixing đến Decision Head; regularization cũng tác động trực tiếp lên các điểm được cộng. Optimizer cập nhật tham số mô hình. Các quyết định truy hồi rời rạc và tài liệu đã được xác định trong Collect. Cơ chế này cho phép khảo sát liệu hành vi truy hồi thích ứng có hình thành mà không cần nhãn thời điểm hay Reinforcement Learning.

H1 khảo sát LM Loss khi thay đổi hệ số trộn \(d\) dưới các tỷ lệ tài liệu hỗ trợ và gây nhiễu. H2 đo gradient theo hệ số được cung cấp trực tiếp để đánh giá hướng của tín hiệu học cục bộ. Hai phép thăm dò kiểm tra giả định nền tảng; đánh giá hệ thống hoàn chỉnh được trình bày riêng ở Chương~\ref{chap:evaluation}.
